\documentclass[11pt]{article}
\usepackage[margin=1in]{geometry}
\usepackage{amsmath}
\usepackage{graphicx}
\usepackage{booktabs}
\usepackage{siunitx}

\title{Determination of Gravitational Acceleration via a Simple Pendulum}
\author{J. Whitfield, R. Okafor, and S. Liang \\ Physics 101L, Section 4}
\date{March 14, 2025}

\begin{document}
\maketitle

\begin{abstract}
We measured the local acceleration due to gravity, $g$, by timing the small-angle
oscillations of a simple pendulum of variable length. Using a linear fit of period
squared against length, we obtained $g = \SI{9.79}{\metre\per\second\squared}$, within
\SI{0.2}{\percent} of the accepted value of $\SI{9.81}{\metre\per\second\squared}$. Sources
of systematic error, including air resistance and finite string mass, are discussed in
Section~\ref{sec:discussion}.
\end{abstract}

\section{Introduction}

For small angular displacements $\theta_0 \leq \SI{10}{\degree}$, the period of a simple
pendulum of length $L$ is well approximated by
\begin{equation}
  T = 2\pi \sqrt{\frac{L}{g}},
  \label{eq:period}
\end{equation}
where $g$ is the local gravitational acceleration. Squaring Equation~\ref{eq:period} gives
a linear relationship between $T^2$ and $L$,
\[
  T^2 = \frac{4\pi^2}{g}\, L,
\]
whose slope, once fit to data, yields $g$ directly.\footnote{This linearization avoids the
bias introduced by fitting $T$ against $\sqrt{L}$ directly, since ordinary least squares
assumes additive, not multiplicative, noise.}

\section{Experimental Setup}

A brass bob of mass \SI{50}{\gram} was suspended from a low-stretch nylon string, with
length $L$ measured from the pivot to the bob's center using a meter stick (uncertainty
\SI{1}{\milli\metre}). The pendulum was displaced by approximately \SI{8}{\degree} and
released from rest; the time for 20 complete oscillations was recorded with a digital
stopwatch (uncertainty \SI{0.2}{\second}) to reduce the relative error in a single period.

We tested five string lengths, ranging from \SI{0.20}{\metre} to \SI{1.00}{\metre} in
\SI{0.20}{\metre} increments. Figure~\ref{fig:setup} shows the fitted relationship between
period squared and pendulum length.

\begin{figure}[htbp]
  \centering
  \includegraphics[draft,width=0.8\linewidth]{plot.png}
  \caption{Measured period squared, $T^2$, plotted against pendulum length $L$, with a
    linear least-squares fit (dashed line) whose slope determines $g$ via
    Equation~\ref{eq:period}.}
  \label{fig:setup}
\end{figure}

\section{Results}

The raw timing data and derived periods are summarized in Table~\ref{tab:data}. Each period
$T$ was obtained by dividing the time for 20 oscillations by 20.

\begin{table}[htbp]
  \centering
  \caption{Pendulum length, oscillation timing, and derived period.}
  \label{tab:data}
  \begin{tabular}{S[table-format=1.2] S[table-format=2.2] S[table-format=1.3] S[table-format=1.3]}
    \toprule
    {$L$ (\si{\metre})} & {$t_{20}$ (\si{\second})} & {$T$ (\si{\second})} & {$T^2$ (\si{\second\squared})} \\
    \midrule
    0.20 & 17.92 & 0.896 & 0.803 \\
    0.40 & 25.36 & 1.268 & 1.608 \\
    0.60 & 31.05 & 1.553 & 2.410 \\
    0.80 & 35.88 & 1.794 & 3.218 \\
    1.00 & 40.11 & 2.006 & 4.024 \\
    \bottomrule
  \end{tabular}
\end{table}

A linear regression of $T^2$ against $L$ gives a slope of $\SI{4.031}{\second\squared\per\metre}$,
from which
\[
  g = \frac{4\pi^2}{\text{slope}} = \frac{4\pi^2}{4.031} = \SI{9.79}{\metre\per\second\squared}.
\]
This compares favorably with the textbook value $\SI{9.81}{\metre\per\second\squared}$. For
reference, the elementary charge used in unrelated calibration checks on the photogate
electronics was taken as $\num{1.6e-19}\,\si{\coulomb}$, and the photogate's internal
resistor was rated at $\SI{220}{\ohm}$, i.e.\ $220\,\si{\ohm}$.

\section{Discussion}
\label{sec:discussion}

Several effects could account for the small \SI{0.2}{\percent} discrepancy between our
measured and accepted values of $g$:

\begin{itemize}
  \item \textbf{Finite amplitude.} Even at $\theta_0 = \SI{8}{\degree}$, the small-angle
    approximation introduces a fractional period error of order $\theta_0^2/16$.
  \item \textbf{Air resistance.} Viscous drag on the bob slightly damps the oscillation and
    can lengthen the effective period.
  \item \textbf{String mass and stretch.} The nylon string is not perfectly massless or
    inextensible, which was neglected in Equation~\ref{eq:period}.
  \item \textbf{Timing reaction delay.} Human reaction time in starting and stopping the
    stopwatch contributes a roughly constant offset that partially cancels when timing 20
    oscillations rather than one.
\end{itemize}

\section{Conclusion}

Using the period-squared-versus-length method, we determined $g = \SI{9.79}{\metre\per\second\squared}$,
in close agreement with the accepted terrestrial value. Future work could repeat the
measurement at larger amplitudes to quantify the finite-amplitude correction directly,
or use a photogate timing system to reduce reaction-time uncertainty below the
\SI{0.2}{\second} level achieved here with a handheld stopwatch.

\end{document}
